% !TEX root = ../journal_0421052099.tex
% ---------------------------------------------------------------------
% The aggregate, and then where it comes from. sec:margin is the point
% of the section: the budget is the second source of a reported gain.
%
% THE GROUNDEDNESS RESULT MUST BE ATTRIBUTED CORRECTLY. Zero ungrounded
% assertions in 1,709 belongs to graph navigation, not to the
% verification layer, and under the grounding definition it is nearly
% guaranteed by construction. Say both, first.
%
% The efficiency result: AGR uses roughly half the CALLS, not half the
% tokens (it spends MORE tokens than ToG, 4,511 vs 3,615 and 6,818 vs
% 5,572). check_paper_numbers.py fails the build on "half the tokens".
%
% tab:togstrata stays in the body: the split conditions on the
% baseline's own outcome, and this table is what makes sec:margin safe
% to read. It is on CWQ and is not on WebQSP; both are said.
%
% THE EQUAL-WIDTH CAVEAT STAYS IN sec:margin. check_paper_numbers takes
% section_body("sec:margin") and requires "clip rate", "read apart" and
% "rather than an aggregate" inside it, because the body proposes the
% re-run and a proposal without its confound reads as a clean
% single-variable experiment, which it is not.
%
% CUT 2026-09-23 (third pass) under a hard 20-22 page budget for
% abstract-to-conclusion, to ~1,150 words from ~2,840. MOVED to
% Appendix D: the wall-clock means, sec:rog, sec:hops, the semantic
% tier with tab:tier2, the equal-width re-run's full argument and the
% entity-breadth check. Each is pointed at from the sentence that used
% to introduce it, and each was chosen because the paper rests no
% conclusion on it. The headline each supports stays here.
% ---------------------------------------------------------------------

\section{Aggregate Results and Where the Margin Lives}\label{sec:results}

The same runs are read several times over: what the five systems answer
correctly and what they spend, then where the margin over the agentic
baseline comes from, and whether the entities they assert are supported
by the graph they navigated. Where AGR is not the best of the five, it
says that too.

\subsection{Accuracy}\label{sec:accuracy}

\Cref{tab:main} gives the main matrix. AGR has the highest Hits@1 and F1
in every column, reaching Hits@1 of $0.755$ on WebQSP against the
agentic baseline's $0.660$ and $0.522$ on ComplexWebQuestions against
$0.435$, so the margin over Think-on-Graph is $9.5$ and $8.7$ points.
The per-question F1 gap is wider still, $0.642$ against $0.477$ and
$0.469$ against $0.378$, because F1 penalises the extra wrong entities
that a set-intersection hit forgives. All eight paired comparisons
against the four baselines reject at the conventional level. The two
against Think-on-Graph, the only comparisons in which both systems
navigate the same graph with the same model under the same call budget,
rest on $76$ questions AGR alone answers against $38$ the other way on
WebQSP ($p = 4.8 \times 10^{-4}$), and $86$ against $51$ on
ComplexWebQuestions ($p = 3.5 \times 10^{-3}$). Against RoG's published
figures, measured over the full splits with a backbone fine-tuned on
both benchmarks, AGR trails on all four; the supplementary material
gives that comparison and what it does and does not establish.

\begin{table}[!tb]
  \begin{center}
    \caption{Main results, $400$ questions per dataset, one frozen
      backbone and one $25$-call budget for every system. Hedge is the
      proportion of questions answered with an empty entity set. The
      three single-call systems are controls that bound a paradigm; the
      comparison the paper reads is the one between the two navigators.}
    \label{tab:main}
    \footnotesize
    % Nine columns is two more than the single-column measure holds at the
    % default 6pt tabcolsep -- it set 2.59pt overfull. The gaps carry no
    % information here, so they give way before the type size does.
    \setlength{\tabcolsep}{4.5pt}
    \begin{tabular}{|l|rrrr|rrrr|}
      \hline
      & \multicolumn{4}{c|}{\textbf{WebQSP}}
      & \multicolumn{4}{c|}{\textbf{ComplexWebQuestions}} \\
      \textbf{System} & Hits@1 & F1 & Hedge & Calls
                      & Hits@1 & F1 & Hedge & Calls \\
      \hline
      \noret        & 0.453 & 0.305 & 12.2\% &  1.0
                    & 0.307 & 0.279 & 38.0\% &  1.0 \\
      \vecrag       & 0.568 & 0.432 & 27.0\% &  1.0
                    & 0.203 & 0.168 & 48.2\% &  1.0 \\
      \graphrag     & 0.340 & 0.262 & 55.8\% &  1.0
                    & 0.205 & 0.183 & 60.8\% &  1.0 \\
      \tog          & 0.660 & 0.477 & 18.8\% & 12.8
                    & 0.435 & 0.378 & 22.5\% & 18.2 \\
      \agr          & \textbf{0.755} & \textbf{0.642} & 8.2\% & 6.2
                    & \textbf{0.522} & \textbf{0.469} & 23.0\% & 8.9 \\
      \hline
    \end{tabular}
  \end{center}
\end{table}

Two features cut against a simple reading. \textbf{Raw hits mislead on
the static graph baseline.} GraphRAG scores $0.340$ on WebQSP, below the
no-retrieval control's $0.453$, which invites the conclusion that it is
broken; the wrong-answer counts invert the picture. The control asserts
on $351$ of $400$ questions and is wrong on $170$ of them, an
\emph{assertion precision} of $51.6\%$, where assertion precision is the
share of questions a system commits on that it answers correctly.
GraphRAG asserts on $177$ and is wrong on $41$, or $76.8\%$, because the
retrieval baselines answer only from the facts retrieved and abstain
when retrieval misses while the control guesses from memory. A grounded
system can score below an unconstrained one on raw hits while committing
four times fewer falsehoods.

\textbf{And AGR does not lead on every column.} On ComplexWebQuestions
its hedge rate of $23.0\%$ is marginally \emph{above} Think-on-Graph's
$22.5\%$, so the accuracy lead is not bought by answering more freely.
The two systems commit at effectively the same rate, on $308$ and $310$
of $400$ questions, and from those near-identical commitments AGR is
right $209$ times and wrong $99$, against $174$ and $136$ for
Think-on-Graph: $67.9\%$ assertion precision against $56.1\%$. It finds
$35$ more answers and asserts $37$ fewer falsehoods than beam search
does, from the same number of commitments.

\subsection{Cost}\label{sec:cost}

AGR reaches those numbers at roughly half the model calls of the agentic
baseline, $6.2$ against $12.8$ on WebQSP and $8.9$ against $18.2$ on
ComplexWebQuestions, ratios of $0.48$ and $0.49$. The call saving is
genuine and the token saving is not: AGR spends $4{,}511$ mean tokens
per question on WebQSP against Think-on-Graph's $3{,}615$, and $6{,}818$
against $5{,}572$ on ComplexWebQuestions. It is the \emph{number of
round trips} that halves, not the volume of text, so on tokens
Think-on-Graph is a frontier point rather than a dominated one, while on
calls it is dominated on both datasets --- and calls are what the shared
budget meters. Wall-clock time tells the same story in seconds (reported
in the supplementary material). AGR never reaches the shared call cap:
across all $800$ test questions it hit the $25$-call limit zero times,
though its depth cap refuses a deeper expansion on $25.6\%$ of test
questions. The baseline reaches the cap far more often.

\subsection{Where the Margin Comes From}\label{sec:margin}

Beam search costs width times depth, so the shared $25$-call cap binds
on Think-on-Graph far more often than on AGR. At width $3$ a full depth
costs the baseline up to thirteen calls (\Cref{sec:baselines}), and the
reimplementation reached its third expansion on $6$ of $400$ WebQSP
questions and $11$ of $400$ on ComplexWebQuestions, so its published
depth of three is nominal under this cap. Splitting the questions by
whether the cap cut it off separates two things the aggregate confounds:
how good its search is, and whether it can afford to finish.
\Cref{tab:togsplit} gives the split, and it reverses the comparison on
one side. Where the budget lets Think-on-Graph run to completion it is
\emph{ahead} of AGR, $0.852$ against $0.788$ on WebQSP and $0.629$
against $0.607$ on ComplexWebQuestions. The entire aggregate margin
comes from the $29\%$ and $44\%$ of questions on which the cap truncates
it, where its accuracy collapses to $0.197$ and $0.188$ while AGR still
reaches $0.675$ and $0.415$.

\begin{table}[!tb]
  \begin{center}
    \caption{AGR against Think-on-Graph, split on whether the shared
      call cap truncated the baseline. Read from the per-record
      \textsf{budget\_exhausted} trace flag.}
    \label{tab:togsplit}
    \footnotesize
    \begin{tabular}{|l|r|rr|}
      \hline
      \textbf{Subset} & \textbf{$n$} & \textbf{\tog{} Hits@1}
        & \textbf{\agr{} Hits@1} \\
      \hline
      \multicolumn{4}{|l|}{\textit{WebQSP}} \\
      \tog{} finished & 283 & \textbf{0.852} & 0.788 \\
      \tog{} clipped  & 117 & 0.197 & \textbf{0.675} \\
      \hline
      \multicolumn{4}{|l|}{\textit{ComplexWebQuestions}} \\
      \tog{} finished & 224 & \textbf{0.629} & 0.607 \\
      \tog{} clipped  & 176 & 0.188 & \textbf{0.415} \\
      \hline
    \end{tabular}
  \end{center}
\end{table}

The split conditions on the baseline's own outcome, so the reading would
be unsafe if the clipped half were simply the hard half, and the
supplementary material reads it within hop strata, and every figure the
reading rests on is quoted here. On WebQSP the clipped half is not
harder for AGR: $81$ of the $117$ clipped questions are single-hop, AGR
scores $0.728$ on them against $0.766$ on the single-hop questions the
baseline finished, and the baseline's fall from $0.846$ to $0.160$
across the same two sets is the cap, not the questions. On
ComplexWebQuestions the halves do differ --- AGR scores $0.309$ on the
clipped single-hop questions against $0.609$ on the finished ones --- so
there truncation and difficulty are confounded and the affordability
reading is weaker. The finished halves are not uniform either: on
ComplexWebQuestions the baseline leads AGR only on single-hop questions,
$0.754$ against $0.609$, and trails on two-hop, $0.561$ against $0.610$,
which is the planner asymmetry of \Cref{sec:planner} seen from the other
side.

The correct reading is affordability, not resolution quality. Exhaustive
beam search finds answers well when allowed to finish; guided search
finds them almost as well and finishes within a budget the beam search
cannot. On a fixed spend the second is the better architecture; the
first would win if spend were unconstrained. The cap of $25$ calls is
itself a choice, and a sweep of it is the first experiment
\Cref{sec:outlook} names.

\textbf{One asymmetry bounds how far that reading extends}, and it is
the candidate-width difference named in \Cref{sec:baselines}. Both
navigators call the same tools but truncate what those tools return at
different widths: Think-on-Graph keeps the first $40$ relations per
entity and the first $20$ neighbours per relation, against AGR's $300$
and $200$. The relation list is sorted by \emph{descending} fanout, so a
binding cut discards the low-fanout tail, which is where the relation
answering a specific question tends to sit. Over the committed tool logs
the $40$-relation cut binds on $31.6\%$ of the $1{,}651$ entities
Think-on-Graph expanded and the $20$-neighbour cut on $32.8\%$ of its
$7{,}992$ neighbour calls, against $1$ of AGR's $3{,}097$ expanded
entities and $3.3\%$ of its neighbour calls.

The affordability conclusion survives this, because a thinner candidate
pool cannot explain why a system runs \emph{out of calls}. What it
bounds is the other half of the split: the $0.852$ and $0.629$ the
baseline reaches where it finishes were measured over a narrower pool
than AGR's, so they are a \textbf{lower bound} on what an equal-width
reimplementation would reach rather than an estimate of it, and since
the baseline is already ahead there, equalising the widths would be
expected to widen that lead. Equalising them is the first change we
would make in a follow-up. That re-run would not vary one thing, though:
wider candidate sets make each pruning call dearer, so equal-width
pruning would also raise the baseline's clip rate and move the boundary
of the very split it would be read from. The two effects have to be read
apart, which means reporting that split at each setting rather than an
aggregate in which they cancel or compound invisibly. The supplementary
material adds the check on whether AGR buys hits by naming more
entities; it does not.

\subsection{Groundedness}\label{sec:groundedness}

Across both datasets AGR asserted $1{,}709$ entities, and \textbf{zero
of them were ungrounded}: every entity it named exists in the
environment and is reachable from a topic entity within the hop cap,
while the parametric control asserted $1{,}001$ entities of which
$22.1\%$ failed that test. Two qualifications come first. The property
is close to guaranteed by construction for any system that copies its
entities out of tool results. And \textbf{it belongs to graph
navigation, not to the verification layer}: Think-on-Graph also reaches
exactly zero ungrounded assertions, on $1{,}265$ entities, without any
claim-checking component whatsoever.

A second tier asks whether an asserted entity is supported \emph{as the
answer}, and it separates the systems far less; the supplementary
material gives every cell, and no between-system gap there survives $n =
60$. The level is the useful finding rather than the ranking.
\textbf{Zero structural hallucination coexists with between a third and
nearly two-thirds of assertions failing a strict semantic-support test},
for every system, including the two that are structurally perfect.
Semantic support is hard for the paradigm, not for one system within it.
